% =====================================================================
% Chapitre 6 — Documentation technique d'exploitation (C2.4.1)
% =====================================================================

\competence{C2.4.1}{Manuel de déploiement, manuel d'utilisation, manuel de mise à jour}{Non}

L'objectif est de permettre à l'équipe CDS d'exploiter et de faire évoluer l'application sans l'intervention de l'auteur. La documentation d'exploitation ne repose pas sur un manuel figé, mais sur trois supports complémentaires et maintenus avec le code~: l'infrastructure versionnée, qui rend le déploiement reproductible et auto-documenté, la documentation d'API générée automatiquement, et le guide d'accompagnement intégré à l'application.

\section{Manuels de déploiement et de mise à jour}
Le déploiement et la montée de version ne s'appuient sur aucune étape manuelle non tracée. Toute la configuration est décrite dans des fichiers versionnés à la racine du dépôt (\texttt{docker-compose.prod.yml}, \texttt{Dockerfile}, \texttt{entrypoint.sh} et \texttt{.gitlab-ci.yml}), ce qui constitue une documentation exécutable et toujours synchronisée avec le code.

Le déploiement initial suppose Docker et Docker Compose installés sur le serveur. Après récupération des variables d'environnement (fichier \texttt{.env} dérivé de \texttt{.env.example}), la commande \texttt{docker compose -f docker-compose.prod.yml up -d} récupère l'image publiée sur le registre GitLab et démarre le conteneur. Au démarrage, le script \texttt{entrypoint.sh} applique automatiquement les migrations Prisma (\texttt{prisma migrate deploy}) avant de servir les requêtes, ce qui évite toute intervention manuelle sur la base.

La mise à jour suit la même logique. Une nouvelle version est livrée sous forme d'image taguée \texttt{vX.Y.Z}, récupérée par \texttt{docker compose pull} puis appliquée par \texttt{up -d}. En cas de problème, le retour arrière consiste à réépingler le tag précédent, sans recompilation. Le détail de cette chaîne, déjà éprouvée jusqu'à la préproduction, est présenté au chapitre~4.

\section{Manuel d'utilisation et documentation de l'API}\label{sec:doc-api}
Côté utilisateur final, la prise en main est assurée directement dans l'application par un guide d'onboarding déclenché à la première connexion, qui accompagne pas à pas la création d'un premier tournoi. Réaffichable à la demande, il tient lieu de manuel d'utilisation intégré. Un manuel pas-à-pas illustré sera complété une fois l'interface stabilisée, les évolutions d'ergonomie issues de la recette étant en cours d'intégration.

Côté développeur, l'API REST est documentée automatiquement par Swagger, généré par NestJS à partir des décorateurs \texttt{@Api*} des contrôleurs. Accessible à l'URL \url{http://localhost:3000/api} en environnement de développement, cette documentation liste chaque endpoint avec ses paramètres et ses schémas de réponse, permet de les tester directement et reste synchronisée avec le code sans maintenance manuelle. La figure ci-dessous en présente l'interface, avec une version pleine page en \hyperref[chap:annexe-swagger]{annexe~\ref*{chap:annexe-swagger}}.

\begin{figure}[H]
    \centering
    \includegraphics[width=0.9\linewidth]{api-docs}
    \caption{Interface Swagger de l'API BrackX (\url{http://localhost:3000/api})~: endpoints regroupés par module (authentification, tournois, administration, tournois publics, utilisateurs) et schémas de données.}
\end{figure}

\livrable{le manuel de déploiement, le manuel d'utilisation et le manuel de mise à jour}
